\section{Proof of the load-axis results}\label{app:loadaxis}
\subsection{Macro moderate deviation probabilities}\label{app:mda}
\begin{proposition}[Moderate deviation probability of a macro core]\label{prop:mda}
Fix $k\ge3$, $d/M=\alpha_c+r_M/\sqrt M$, with $r_M\to-\infty$ and
$|r_M|\le C\sqrt{\log M}$. Let $\mathcal A$ be the event that the actual 2-core has more than
$\varepsilon_0M$ edges, with $0<\varepsilon_0<\alpha_cq_c^k$. Then in the Plain model
$\Pr(\mathcal A)=\Phi(z_r)(1+o(1))$, where $z_r=-\beta_kr_M/\sigma_k$.
\end{proposition}
Proposition~\ref{prop:mda} gives the relative probability of a macro failure between the critical
window and low load. Equation (2.18) of Dembo--Montanari gives an expansion for the configuration
model for each fixed $r$, with additive error of order $n^{-\eta}$ ($\eta<5/26$)\cite{dm}. The
additive error converts to relative precision only when $\Phi(-t)\gg n^{-\eta}$.
Proposition~\ref{prop:mda} gives a relative asymptotic over the entire range
$|r_M|=O(\sqrt{\log M})$ and transfers the result to the Plain model.

The next four subsections prove Proposition~\ref{prop:mda}. Write $L=\lceil\log M\rceil$,
$\vartheta_M=\sqrt M(\log M)^6$, $\varepsilon_M=M^{-1/8}$, and $t=-z_r$. ``Superpolynomially
small'' means $o(M^{-b})$ for every fixed $b>0$.

\subsubsection{The event on which the concentration estimates hold}
The estimates used for the uniform concentration before stopping in Appendix~\ref{sec:main-proof}
hold uniformly on a fixed compact neighborhood of $\mu_M=kd/M$. Concretely: the tail probability
of the maximum degree is superpolynomially small; for a fixed target $(n,j)$ the difference of the
backward initial-state functional has a constant upper bound, the dynamic martingale difference
also has a constant upper bound, and with threshold $\sqrt M(\log M)^2$ the McDiarmid and Azuma
inequalities give a tail of $\exp(-c\log^4M)$; under the maximum-degree truncation the accumulated
drift is $O(L^4)$. Taking a union over $O(ML)$ targets yields a good event $G_M$ whose complement
has superpolynomially small probability. On $G_M$, up to the stopping time (including the stopping
instant itself), the error of every histogram is at most $CB_M$. Before $n_-$ the fluid degree-one
density is at least $cM\varepsilon_M^2-C\sqrt M|r_M|$, which is far larger than $\vartheta_M$, so
on $G_M$ the process does not stop before $n_-$.

\subsubsection{Concatenating the initial and dynamic stages into a martingale}
The configuration model can be generated as follows: first draw $kd$ iid uniform endpoints, then
assign the labeled stubs to the slots of the edges independently and uniformly. The initial-state
filtration records the endpoint sequence; the dynamic filtration adds to this the part of the
pairing revealed so far by the peeling process. Given the endpoint sequence, the unrevealed pairing
is still uniform, so the initial Doob increments and the dynamic increments can be concatenated
into a single martingale difference sequence. For $h_c(j)=j(1-a_c)^{j-1}$, the unnormalized Doob
increment of the $(l+1)$-th endpoint is
\begin{gather*}
\Delta_l=\nu_{n-l-1}\bigl(D_V^{(l)}\bigr)-\frac1M\sum_v\nu_{n-l-1}\bigl(D_v^{(l)}\bigr),\\
\nu_m(j)=\E h_c\bigl(j+1+\mathrm{Bin}(m,1/M)\bigr)-\E
h_c\bigl(j+\mathrm{Bin}(m,1/M)\bigr),
\end{gather*}
This increment is bounded. Changing one endpoint changes the initial-state bracket $V_{\rm in}$ by
at most $C/M$, so by McDiarmid's inequality $V_{\rm in}$ concentrates within a band of width
$(\log M)^2/\sqrt M$. A direct computation of $\Var\left(\sum_vh_c(D_v)\right)$ (retaining the
$1/M$ terms in the two-vertex covariance) gives
$\E V_{\rm in}=\sigma_0^2/\sigma_k^2+O((1+|r_M|)/\sqrt M)$. On $G_M$ the dynamic bracket differs
from $\tau_k^2/\sigma_k^2$ by $O(M^{-1/8}+(\log M)^8/\sqrt M+|r_M|/\sqrt M)$. Hence the
concatenated martingale $S_M$ has increments of size at most $C/\sqrt M$; its bracket has a
deterministic upper bound and, except for a superpolynomially small exceptional probability, lies
within $1\pm v_M$, where $v_M=C\{M^{-1/8}+(\log M)^8/\sqrt M\}$.

\subsubsection{Exponential tilting and the relative tail probability}
Let $m_j=\E[e^{-t\xi_j}\mid\mathcal F_{j-1}]$ and define the measure $Q$ with density
$\prod_je^{-t\xi_j}/m_j$. On the good set,
$\prod_jm_j=e^{t^2/2}(1+O(t^2v_M+t^3/\sqrt M))$. The bound $Z_t^2\le Z_{2t}e^{Ct^2}$ gives the
second-moment bound $\E_P(dQ/dP)^2\le e^{Ct^2}$, so the bad set remains negligible under $Q$ and
the total variation between the exponential tilting $P_t$ and $Q$ tends to zero. Under $Q$ the
conditional mean of the increment is $-tv_j+O(t^2v_j/\sqrt M)$; recentering then makes the bracket
close to $1$. For a centered martingale with increment bound $\beta'$ and bracket exceptional
probability $q'$, the telescoping expansion of the compensating exponential and Esseen smoothing
give
\begin{equation*}
\sup_x|P(T\le x)-\Phi(x)|\le C(\beta'+v'+q')^{1/8}+Cq'.
\end{equation*}
At the magnitudes of the present model, the Kolmogorov error $\epsilon_M$ of the tilted
distribution satisfies $t\epsilon_M\to0$. The untilting formula is
\begin{equation*}
\Pr(S_M\le-t)=\E e^{-tS_M}\,e^{-t^2}\,\E_t\bigl[e^{t(S_M+t)}\mathbf1_{\{S_M+t\le0\}}\bigr]
\end{equation*}
Integrating by parts gives $\Pr(S_M\le-t)=\Phi(-t)(1+o(1))$. The same argument gives
$\Pr(|S_M+t|\le\delta)=o(\Phi(-t))$ for bandwidths $\delta$ satisfying $t\delta\to0$.

\subsubsection{Transfer to the actual macro failure and the Plain model}
On $G_M$, $R_M/\sigma_k=S_M+b_M+O((\log M)^4/\sqrt M)$; the exact mean expansion gives
$b_M=-g_kr_M+O((1+r_M^2)/\sqrt M)$, where $g_k=ka_ce^{-x_c}/\sigma_k$. Replacing the $O_p$ bound
from the local freezing by a deterministic bandwidth $\delta_M\to0$ (satisfying
$b_M\delta_M\to0$) at the cost of a superpolynomially small exceptional probability gives
\begin{equation*}
\mathcal A\,\triangle\,\{S_M\le-b_M\}\subseteq\mathcal B_M\cup\{|S_M+b_M|\le C\delta_M\},
\end{equation*}
where $\mathcal B_M$ has superpolynomially small probability. When the score lies below the
negative bandwidth the process stops while the residual is still of linear size, so it lies in
$\mathcal A$; when the score lies above the positive bandwidth the process does not stop in the
window or on either side of it, until the residual drops below $\varepsilon_0M/2$, so it does not
lie in $\mathcal A$. This step does not require excluding small cores. Combining the relative tail
estimate with the bandwidth estimate gives the result in the configuration model.

Then transfer to the Plain model. When $l$ special edges are fixed, the change of
$\sum_vh_c(D_v)$ induced by the special edges has a deterministic upper bound and the per-step
error of the sampling kernel is $O_l(1/M)$, so all the rates above remain unchanged and
$\Pr(\mathcal A\mid I_1=\cdots=I_l=1)/p_M\to1$. Applying finite-order upper and lower Bonferroni
truncations to the indicator of having no bad edges, first letting $M\to\infty$ and then letting
the truncation order tend to infinity, gives $\Pr(\mathcal A,N^{\rm bad}_M=0)/p_M\to
e^{-\Lambda_k}$; dividing by $\Pr(N^{\rm bad}_M=0)\to e^{-\Lambda_k}$ gives the result in the
Plain model. The key point is that every fixed planted term already has precision relative to
$p_M$.
\subsection{Fixed subcritical load}\label{app:lowload}
\begin{proposition}[Fixed subcritical load]\label{prop:lowload}
Fix $k\ge3$ and $d/M\to\alpha\in(0,\alpha_c)$. Write $\lambda=\binom d2/\binom Mk$, let $P_2$ be
the event that the actual core is exactly two repeated edges, and let
$\gamma=(d\rho_M^2-1)/(d-1)$ be the sign factor of Section~\ref{sec:sign}. There exists
$\delta_k\ge1$ such that
\begin{equation*}
\Pr(F)=\lambda\{1+O(M^{-\delta_k})\},\qquad \Pr(P_2\mid F)\to1,\qquad X\mid F\Rightarrow N(0,1),
\end{equation*}
\begin{equation*}
\E[X\mid F]=\frac{2\gamma}{s_M}+O(|\gamma|M^{-\delta_k}),\qquad c_F\binom d2\longrightarrow1.
\end{equation*}
\end{proposition}
Proposition~\ref{prop:lowload} shows that at a fixed subcritical load the failure is almost always
caused by two repeated edges; in this regime the standardized energy tends again to normality, and
the conditional mean on the original scale retains the leading term $2\gamma/s_M$, whose origin is
explained in Section~\ref{sec:loadaxis}.

This subsection proves Proposition~\ref{prop:lowload}. Write $K=\binom Mk$, $J=\binom d2$,
$N=\sum_{i<j}\mathbf1_{\{S_i=S_j\}}$, $D=\{N\ge1\}$, and $\lambda=J/K$. Let $a=(k-2)/2$,
$b_k=\lceil3(k-2)/2\rceil$, $\delta_k=\min\{k-2,\,b_k-(k-2)\}\ge1$,
$W_{ij}=|S_i\cap S_j|-k^2/M$, and
$Y=\widehat d-d=\frac{2}{k(1-k/M)}\sum_{i<j}s_is_jW_{ij}$.

\subsubsection{Uniform counting of small cores}
A 2-core subgraph with $e$ edges uses at most $v_e=\lfloor ke/2\rfloor$ vertices. For a prescribed
set $I$ of $e$ edge labels and a $v_e$-element vertex set $V$, let
$A_{I,V}=\{S_i\subseteq V,\ i\in I\}$; every such core lies in some $A_{I,V}$. The first moment
is
\begin{equation*}
B_e=\binom de\binom M{v_e}\Bigl(\binom{v_e}k\Big/\binom Mk\Bigr)^e\le\bigl[C(e/M)^a\bigr]^e,
\end{equation*}
This bound holds uniformly for $d\le AM$ and $e\le\varepsilon M$ ($\varepsilon<1/k$). For $e=3$,
retaining the rounding gives $B_3=O(M^{-b_k})$. For $e\ge4$ the ratio of adjacent terms is at
most $Ce^a((e+1)/M)^a$. For fixed $r\ge0$, choose $\varepsilon$ small enough that the adjacent
ratio of $e^rB_e$ is at most $1/2$ on $4\le e<\varepsilon M$; then
\begin{equation}
\sum_{e=3}^{\lfloor\varepsilon M\rfloor}e^rB_e=O(M^{-b_k}).\label{eq:smallcore}
\end{equation}
This sum also covers the intermediate range with $e\to\infty$ and $e=o(M)$.

\subsubsection{Superpolynomially small probability of a linear-size core}
For every fixed $\varepsilon>0$ and $B>0$, $\Pr(|\mathrm{core}|\ge\varepsilon M)=o(M^{-B})$. The
proof uses message passing of fixed depth. Let $Z_T$ be the number of edges whose incoming
messages are all $1$ at round $T$; then $|\mathrm{core}|\le Z_T$ on every graph. Local convergence
to a Poisson branching tree gives $\E Z_T/M\to\alpha q_T^k$, where $q_0=1$ and
$q_{t+1}=1-\exp(-k\alpha q_t^{k-1})$. When $\alpha<\alpha_c$ this recursion has no positive fixed
point, so one can choose $T$ with $\alpha q_T^k<\varepsilon/4$. On the set $H$ of graphs with
maximum degree at most $\lceil\log M\rceil$, changing $h$ edge coordinates changes $Z_T$ by at
most $L_Mh$, where $L_M$ is of polylogarithmic order. Applying a McShane extension of $Z_T|_H$
truncated to $[0,d]$, and then McDiarmid's inequality, gives
$\Pr(Z_T\ge\varepsilon M)\le\Pr(H^c)+\exp\{-cM/L_M^2\}=o(M^{-B})$.

\subsubsection{Failure probability and the conditional law}
A single edge cannot form a core; two edges form a core if and only if they are identical. Hence
on $F\setminus P_2$ the actual core has at least three edges, and by the previous two subsections
$\Pr(F\setminus P_2)=O(M^{-b_k})$. The probability that any two distinct edge pairs are
simultaneously repeated is $K^{-2}$, so $\E N=\lambda$ and $\E[N(N-1)]\le\lambda^2$; by
Bonferroni's inequality $\Pr(F)=\lambda+O(\lambda^2+M^{-b_k})$.

Let $Q$ be the distribution obtained by planting a repeated pair at random, whose density with
respect to the original distribution is exactly $N/\lambda$; then
$\|P(\cdot\mid D)-Q\|_{\rm TV}=O(\lambda)$. Since $D\subseteq F$,
$\|P(\cdot\mid F)-P(\cdot\mid D)\|_{\rm TV}=\Pr(F\setminus D)/\Pr(F)$. Combining the two bounds
gives $\|P(\cdot\mid F)-Q\|_{\rm TV}=O(M^{-\delta_k})$. This is a total variation bound over the
entire labeled mapping graph.

\subsubsection{CLT under the unconditional law and under the repeated-edge conditioning}
Let $v_i=\mathbf1_{S_i}-(k/M)\mathbf1$ and reveal the edges in order of their labels. The
martingale difference is $\Delta_j=\frac{2}{k(1-k/M)}s_jv_j^\top\sum_{i<j}s_iv_i$, and
$Y=\sum_j\Delta_j$. The sum of conditional variances is
\begin{gather*}
\mathcal
V_d=s_M^2+\frac{8}{k(1-k/M)(M-1)}\sum_{i<\ell}(d-\ell)s_is_\ell
W_{i\ell},\\
\Var\mathcal
V_d=\frac{64}{(M-1)^3}\sum_{\ell=2}^d(\ell-1)(d-\ell)^2=O(M),
\end{gather*}
so $\mathcal V_d/s_M^2$ converges to $1$ in $L^2$. The bound $\E\Delta_j^4=O(1)$ gives the
conditional Lindeberg condition, so the martingale CLT yields $X\Rightarrow N(0,1)$, uniformly in
any deterministic signs. Fix a planted pair $(a,b)$, write their common mapping as $T$ and their
sign product as $\sigma$; then
\begin{equation*}
Y=Y_R+2\sigma+\frac{2}{k(1-k/M)}(s_a+s_b)\sum_{j\in R}s_j\bigl(|T\cap S_j|-k^2/M\bigr),
\end{equation*}
The variance of the last term is $O(1)$. Applying the unconditional CLT to the remaining edges and
then transferring back to $P(\cdot\mid F)$ via the total variation bound completes the proof.

\subsubsection{Expansion of the conditional mean}
The total variation bound does not by itself control the mean, so
$\E[|Y|\,|\mathbf1_F-N|]$ must be estimated directly. Given $A_{I,V}$ (with $|I|=e$), splitting
$Y$ into an internal, a cross, and an external part gives $\E[|Y|\mid A_{I,V}]\le C(\sqrt
M+e^2)$. By \eqref{eq:smallcore} (with $r=0,2$) and the superpolynomial bound for linear-size
cores, $\E[|Y|;F,N=0]=O(\sqrt M M^{-b_k})$. When two prescribed repeated pairs occur
simultaneously, $\E|Y|\le C\sqrt M$, so $\E[|Y|(N-1);N\ge1]\le C\sqrt M\lambda^2$. When a
prescribed edge pair $(i,j)$ is repeated, $\E Y=2s_is_j$, so $\E[YN]=2\gamma\lambda$, and hence
\begin{equation*}
\E[Y;F]=2\gamma\lambda+O\bigl(\sqrt M(M^{-b_k}+\lambda^2)\bigr).
\end{equation*}
Since $F$ is invariant under permutations of the edge labels, for each $M$ exactly
$\E[X\mid F]=\gamma\E[X_+\mid F]$, where $X_+$ is the measurement with all-positive signs.
Applying the above identity in the all-positive case and multiplying by $\gamma$ gives the
remainder with factor $|\gamma|$ in the mean expansion. Matching the coefficients with
\cite{prior} in the all-positive case gives
$c_F\binom d2=\tfrac12\E[Y_+\mid F]=1+O(M^{1/2-\delta_k})$.
\subsection{The macro conditional law and the two-mechanism mixture}\label{app:mixture}
\subsubsection{Setup and the results used}

This subsection fixes $k\ge3$, uses the Plain model, and places the load in the moderate
deviation regime:
\begin{equation}
\frac dM=\alpha_c+\frac{r_M}{\sqrt M},\qquad r_M\to-\infty,\qquad |r_M|\le C\sqrt{\log M}.\label{F:range}
\end{equation}
Write $t_M=-g_kr_M\to\infty$, where $g_k=ka_ce^{-x_c}/\sigma_k$;
$p_M=\Phi(-t_M)$; and $\lambda_M=\binom d2/\binom Mk$. The macro event is
$\mathcal A=\{\text{number of edges of the actual 2-core}>\varepsilon_0M\}$. The sign factor is
$\gamma=(d\rho_M^2-1)/(d-1)$; the all-positive measurement is denoted $X_+$, and the measurement
under an arbitrary deterministic sign composition is denoted $X_s$. The left-tail Mills ratio is
written as $A(-t)=\phi(t)/\Phi(-t)$ following the main text. Let
\begin{equation*}
G_{M,s}(x)=\frac{\Phi_2(x,-t_M;\eta_k\rho_M^2)}{\Phi(-t_M)}.
\end{equation*}

\begin{proposition}[Macro conditional law and two-mechanism mixture]\label{prop:mixture}
Within the range \eqref{F:range}, take $\varepsilon_0$ small enough (Lemma~\ref{F:lem:cutoff}),
let the signs be an arbitrary deterministic composition, and do not require $\rho_M$ to converge.
Then
\begin{gather*}
\sup_x\bigl|\Pr(X\le x\mid \mathcal A)-G_{M,s}(x)\bigr|\to0,\qquad
\Var(X\mid \mathcal A)=1-\eta_k^2\rho_M^4+o(1),\\
\sup_x\left|\Pr(X\le x\mid
F)-\frac{p_MG_{M,s}(x)+\lambda_M\Phi(x)}{p_M+\lambda_M}\right|\to0.
\end{gather*}
\end{proposition}
The first two displays say that in the moderate deviation regime a macro failure still obeys the
selection law of the critical window, with only its center shifted by $t_M$. The last display says
that the failure-conditioned distribution is a two-layer weighted mixture whose weights $p_M$ and
$\lambda_M$ are both explicit.

This subsection and Appendix~\ref{app:wm} prove four groups of statements, which correspond to
the main text as follows.
\begin{enumerate}[label=(\arabic*)]
\item The macro conditional CDF with all-positive signs:
$\sup_x|\Pr(X_+\le x\mid \mathcal A)-\Phi_2(x,-t_M;\eta_k)/\Phi(-t_M)|\to0$ (Proposition~\ref{F:prop:cdf}).
\item The conditional moments with all-positive signs and the conditional mean under an arbitrary
composition (Proposition~\ref{F:prop:moments}).
\item The macro conditional CDF and conditional variance under an arbitrary composition, namely
the first two displays of Proposition~\ref{prop:mixture} (Proposition~\ref{F:prop:sign}).
\item The control of the three weighted error terms, namely the weighted mean formula of
Proposition~\ref{prop:wm}, proved in Appendix~\ref{app:wm}.
\end{enumerate}

We use directly the scalar inputs already established in Appendix~\ref{app:mda} (the proof of
Proposition~\ref{prop:mda}). For ease of reference they are collected here.

\begin{enumerate}[label=(C\arabic*)]
\item There exists a good event $G_M$ whose complement has superpolynomially small probability. On
$G_M$, up to the stopping time, the error of every histogram is at most $C\vartheta_M$, with
$\vartheta_M=\sqrt M(\log M)^6$.
\item In the configuration model, the initial Doob increments and the dynamic increments
concatenate into a score martingale $S$; its increments are at most $C/\sqrt M$, its bracket has a
deterministic upper bound, and except for a superpolynomially small exceptional probability it
lies within $1\pm v_M$, with $v_M=C\{M^{-1/8}+(\log M)^8/\sqrt M\}$.
\item There is a bandwidth sandwich between the macro event and the score event:
$\mathcal A\,\triangle\,\{S\le-b_M\}\subseteq\mathcal B_M\cup\{|S+b_M|\le C\delta_M\}$, where
$\mathcal B_M$ is superpolynomially small, the deterministic bandwidth $\delta_M\to0$ satisfies
$b_M\delta_M\to0$, and $b_M=t_M+O((1+r_M^2)/\sqrt M)$.
\item The relative asymptotic of the macro probability: $\Pr(\mathcal A)=p_M(1+o(1))$; the same
holds with $l$ fixed special edges, and transfers to the Plain model via a relative Bonferroni
bound.
\end{enumerate}

``Superpolynomially small'' means $o(M^{-b})$ for every fixed $b>0$. By \eqref{F:range}, $p_M$ is
no smaller than a fixed negative power divided by a polylogarithm, so a superpolynomially small
quantity divided by $p_M$ is still superpolynomially small. This is used repeatedly below.

\subsubsection{Choosing the lower bound on the macro core edge count}

Proposition~\ref{prop:mda} holds for any $0<\varepsilon_0<\alpha_cq_c^k$. The two-mechanism
reduction also needs the small-core counting to cover all core sizes below the macro layer, so
here a small enough cutoff is fixed.

\begin{lemma}[Choice of the cutoff]\label{F:lem:cutoff}
Take $\varepsilon_*>0$ such that the small-core sum \eqref{eq:smallcore} holds for
$r\in\{0,2\}$ on $3\le e\le\varepsilon_*M$, and $\varepsilon_*<\alpha_cq_c^k/2$. Fix
$0<\varepsilon_0\le\varepsilon_*$. Let $D$ be the event that some repeated edge pair exists, and
$R=F\setminus(\mathcal A\cup D)$. Then under \eqref{F:range}
\begin{equation*}
\Pr(R)=O(M^{-b_k})=o(\lambda_M),\qquad \Pr(\mathcal A\cap D)=o(\lambda_M).
\end{equation*}
\end{lemma}

\begin{proof}
On $R$ the actual core has between $3$ and $\varepsilon_0M$ edges. The first-moment sum of
Appendix~\ref{app:lowload} requires only $d\le AM$, independent of the distance to the threshold,
so $\Pr(R)=O(M^{-b_k})$; and since $b_k-(k-2)\ge1$, this is $o(\lambda_M)$.

For the intersection, first consider a fixed window. Let
$d/M=\alpha_c+r/\sqrt M+o(M^{-1/2})$ with $r$ fixed, plant a pair of identical uniform
$k$-subsets, and let the remaining edges be independent. The finite-planting lemma of
Appendix~\ref{sec:mp-plain} allows the special edges to be correlated, so the precritical score
still has the same Gaussian limit; the two repeated edges are themselves valid $k$-subsets, the
probability of no bad edge among the remaining ordinary edges has a positive limit, and the
inclusion-exclusion transfer proceeds as usual. The planted repeated edge forces the final
failure, so here only the macro event is transferred: when the score lies below a fixed negative
gap the process stops with $(\alpha_cq_c^k+o(1))M$ edges remaining, which lies in $\mathcal A$;
when the score lies above a fixed positive gap the pre-stopping concentration and local freezing
rule out early stopping, until the residual drops below $\varepsilon_0M/2$, which does not lie in
$\mathcal A$. Removing the safety gap gives $Q_r(\mathcal A)\to\Phi(g_kr)$, where $Q_r$ is the
planted distribution.

Return to \eqref{F:range}. For any fixed $r_0<0$, eventually
$d\le\lfloor\alpha_cM+r_0\sqrt M\rfloor$. Coupling by adding the same sequence of ordinary edges,
the macro event is monotone under edge addition, so
$\limsup_MQ_M(\mathcal A)\le\Phi(g_kr_0)$. Letting $r_0\to-\infty$ gives $Q_M(\mathcal A)\to0$.
The density of the random-pair planted distribution with respect to the original distribution is
$N/\lambda_M$, where $N$ is the number of repeated edge pairs, so
\[
\Pr(\mathcal A\cap D)\le\E[N\mathbf1_{\mathcal A}]=\lambda_MQ_M(\mathcal A)=o(\lambda_M).\qedhere
\]
\end{proof}

The probability asymptotic of Proposition~\ref{prop:mda} holds for any
$\varepsilon_0<\alpha_cq_c^k$, so the macro probability itself is independent of the choice of
cutoff. Proposition~\ref{prop:mixture}, Proposition~\ref{prop:wm}, and the results in this
subsection conditioned on $\mathcal A$ all refer to the cutoff selected in
Lemma~\ref{F:lem:cutoff}.

\subsubsection{Initial-state Doob decomposition of the measurement and the score with all-positive signs}

Work first in the configuration model. Let $n_e=kd$ be the total number of endpoints,
$\mu=n_e/M$, the endpoints $V_1,\ldots,V_{n_e}$ independent and uniform, and $D_v$ the occupancy
count. Define
\begin{equation*}
K=\sum_v(D_v)_2,\qquad Y_K=\frac{K-\E K}{\sqrt2\mu\sqrt M},\qquad H=\sum_vh_c(D_v),
\end{equation*}
where $(j)_2=j(j-1)$, $h_c(j)=j(1-a_c)^{j-1}$, and $\E K=n_e(n_e-1)/M$. $Y_K$ is the
standardized occupancy energy, which is converted back to $X_+$ in \eqref{F:affine} at the end.

After revealing the first $l$ endpoints, write the occupancy count as $D_v^{(l)}$ and let
$m=n_e-l-1$. For a function $f$ define
$\nu^f_m(j)=\E[\Delta f(j+\mathrm{Bin}(m,1/M))]$, where $\Delta f(j)=f(j+1)-f(j)$. When the next
endpoint $V$ falls, the Doob increment of $\sum_vf(D_v)$ is
\begin{equation*}
\Delta_f=\nu^f_m(D_V^{(l)})-\frac1M\sum_v\nu^f_m(D_v^{(l)}).
\end{equation*}
Taking $f(j)=(j)_2$ gives $\nu^f_m(j)=2j+2m/M$, so the Doob increment of the energy is exactly
$2(D_V^{(l)}-l/M)$. The total number of endpoints is fixed sample by sample; here the Poisson
ensemble does not replace the true initial state.

\begin{lemma}[Mean of the bracket]\label{F:lem:bracket}
Let the matrix bracket of the initial-state vector $(Y_K,(H-\E H)/(\mu_c\sigma_k\sqrt M))$ be
$(V_{EE},V_{EH},V_{HH})$. Then
\begin{equation}
\E V_{EE}=(1-1/n_e)(1-1/M),\qquad \E V_{EH}=\eta_k+O\bigl((1+|r_M|)/\sqrt M\bigr),\label{F:bracketmean}
\end{equation}
$\E V_{HH}$ is the same as the initial-state bracket of Appendix~\ref{app:mda}.
\end{lemma}

\begin{proof}
Expanding the collision indicators gives $\Var K=2n_e(n_e-1)(M-1)/M^2$. For any $h$, let
$B\sim\mathrm{Bin}(n_e-2,1/M)$; there is the finite identity
\begin{equation*}
\Cov(K,H)=\frac{n_e(n_e-1)(M-1)}{M^2}\E\Delta^2h(B).
\end{equation*}
Here is the proof. Write $K=\sum_{i\ne j}\mathbf1_{\{V_i=V_j\}}$. Conditioned on $V_1=V_2$, the
mean of $H$ is $\E h(B+2)+(M-1)\E h(B)$; the unconditional mean is
$M\E h(B)+2\E\Delta h(B)+M^{-1}\E\Delta^2h(B)$. Their difference is
$(1-M^{-1})\E\Delta^2h(B)$; multiplying by $n_e(n_e-1)/M$ gives the claim.

For $h=h_c$, the Poisson limit gives
$\E_{\Pois(\mu)}\Delta^2h_c(D)=a_ce^{-\mu a_c}(\mu a_c-2)$. The function $\Delta^2h_c$ is
bounded, the total variation error between Binomial and Poisson is $O(M^{-1})$, and the parameter
$\mu$ lies in a compact interval and the dependence is smooth. By the Doob isometry, $\E V_{EE}$
and $\E V_{EH}$ are the corresponding variance and covariance divided by the normalization
constant; substituting $\mu=\mu_c+O(|r_M|/\sqrt M)$ gives \eqref{F:bracketmean}. The limit of the
cross term is exactly $\eta_k$, which comes from the true occupancy model and does not presuppose
independence between the dynamic noise and the energy.
\end{proof}

The maximum occupancy count can be of order $\log M$, so the energy increments are unbounded. To
obtain concentration, take $L=\lceil\log M\rceil$ and define the truncated energy $e_L$ by
$e_L(0)=0$ and $\Delta e_L(j)=2\min(j,L)$. For $j\le L+1$ we have $e_L(j)=(j)_2$, and
$|\Delta e_L|\le2L$, $|\Delta^2e_L|\le2$. The true Doob martingale is reconstructed from $e_L$;
the large increments of the original martingale are not simply deleted.

\begin{lemma}[Superpolynomial concentration of the truncated bracket]\label{F:lem:conc}
For $(f,g)\in\{(e_L,e_L),(e_L,h_c),(h_c,h_c)\}$,
\begin{equation}
\Pr\Bigl(|V^{(L)}_{fg}-\E V^{(L)}_{fg}|>CL^2(\log M)^2/\sqrt M\Bigr)\le2e^{-c(\log M)^4},\label{F:conc}
\end{equation}
and the difference between $\E V^{(L)}_{fg}$ and the mean of the untruncated bracket, multiplied
by any fixed power of $M$, still tends to zero.
\end{lemma}

\begin{proof}
Changing one input endpoint changes, at each subsequent $l$, two occupancy coordinates by $1$
each. For $e_L$, $\nu^{e_L}_m$ is bounded by $O(L)$ and a single coordinate changes by $O(1)$;
for $h_c$ both are $O(1)$. So each
$q_l^{fg}=M^{-1}\sum_v\nu^f_m\nu^g_m-\bar\nu^f_m\bar\nu^g_m$ changes by at most $CL^2/M$; after
normalization and summing over $O(M)$ steps, the sensitivity of the whole bracket to each input
coordinate is at most $CL^2/M$. McDiarmid's inequality gives \eqref{F:conc}.

The probability that the maximum occupancy count exceeds $L$ is at most $M(eC/L)^L$, smaller than
any fixed negative power of $M$. On its complement $K$ equals the truncated energy, and both are
at most $O(M^2)$ in absolute value on the whole sample space. Hence the differences in the center,
the variance, and the covariance with $H$, multiplied by any fixed power of $M$, still tend to
zero. Here only the replacement of the mean is needed; the truncated bracket need not equal the
original bracket path by path.
\end{proof}

\subsubsection{The dynamic stage and the deterministic upper bound on the bracket}

After the initial state, the labeled stubs are assigned to the edges independently and uniformly,
and the dynamic filtration follows the construction of Appendix~\ref{app:mda}, without revealing
the full pairing in advance. The initial-state increment of the two-dimensional vector is
$(\Delta e_L/(\sqrt2\mu\sqrt M),\Delta h_c/(\mu_c\sigma_k\sqrt M))$, and the dynamic increment is
$(0,\text{score martingale increment})$. In the dynamic stage the first coordinate is zero, so it
contributes no new $EE$ or $EH$ bracket; this is an identity on the filtration and does not
require an independence assumption.

By Lemma~\ref{F:lem:bracket}, Lemma~\ref{F:lem:conc}, and input (C2), except for a
superpolynomially small exceptional set the total matrix bracket $\mathcal V$ satisfies
\begin{equation}
\|\mathcal V-\Sigma\|\le v_M,\qquad \Sigma=\begin{pmatrix}1&\eta_k\\\eta_k&1\end{pmatrix}.\label{F:Sigma}
\end{equation}
$\Sigma$ is positive definite: the initial-state energy has variance $1$, and the score has
strictly positive dynamic variance $\tau_k^2$, so $|\eta_k|<1$.

After truncation the increment bound is $\beta_M=CL/\sqrt M$, but on the whole sample space the
trace of the bracket is only known to be $O(L^2)$. To perform the exponential tilting, the entire
vector is stopped when the trace of the bracket first exceeds the fixed constant
$C_0>\operatorname{tr}\Sigma+1$. The bracket increments are positive semidefinite matrices, so the
trace is monotone; the stopping changes the terminal value only when the total trace exceeds
$C_0$, an event of superpolynomially small probability. After the stopping,
$\operatorname{tr}\mathcal V\le C_0+C\beta_M^2$ holds deterministically. Below, $(Y^*,S^*)$
denotes this stopped vector. On the good set it equals the original score and energy up to a
deterministic center term that is superpolynomially small. The subsequent exponential tilting acts
only on $S^*$.

\subsubsection{Error rate of the two-dimensional distribution function under exponential tilting}

This gives the explicit polynomial rate needed for the macro conditional law.

\begin{lemma}[Single-coordinate tilting lemma]\label{F:lem:tilt}
Let a two-dimensional martingale difference $\zeta_j=(u_j,v_j)$ satisfy $|\zeta_j|\le\beta$, let
the trace of the total matrix bracket $\mathcal V$ have a deterministic upper bound, and let
$\Pr(\|\mathcal V-\Sigma\|>v)\le q$, where $\beta,v$ are polynomially small in $M$ and $q$ is
superpolynomially small. Write $(U,S)=\sum_j\zeta_j$, let $1\le t\le C\sqrt{\log M}$, and define
the tilted measure $dP_t/dP=e^{-tS}/\E e^{-tS}$. Then the lower-rectangle CDF distance
$\epsilon_M$ between $(U+\eta_kt,S+t)$ under $P_t$ and $N(0,\Sigma)$ satisfies
\begin{equation}
\epsilon_M\le C\Bigl\{(\beta+v+t\beta+q_Q)^{1/64}+tv+t^2\beta+t^2v+t^3\beta+q_Q\Bigr\},\label{F:rate}
\end{equation}
where $q_Q$ is superpolynomially small. Substituting $\beta=O(\log M/\sqrt M)$ and
\eqref{F:Sigma}, the right-hand side is a fixed negative power times a polylogarithm, so
$t\epsilon_M\to0$.
\end{lemma}

\begin{proof}
\emph{Step 1 (product tilting).} Define the measure $Q$ with density
$\prod_je^{-tv_j}/\E(e^{-tv_j}\mid\mathcal F_{j-1})$. The deterministic trace bound gives
$\E(dQ/dP)^2\le e^{Ct^2}$, so the bad set remains superpolynomially small under $Q$; the total
variation distance between the true tilting $P_t$ and $Q$ is $O(t^2v+t^3\beta)$ plus a
superpolynomially small term.

\emph{Step 2 (drift and bracket).} Write
$\mathcal B_j=\E(\zeta_j\zeta_j^\top\mid\mathcal F_{j-1})$ and $e_2=(0,1)^\top$. A stepwise
Taylor expansion gives
\begin{equation}
\E_Q(\zeta_j\mid\mathcal F_{j-1})=-t\mathcal B_je_2+\varrho_j,\qquad|\varrho_j|\le Ct^2\beta\operatorname{tr}\mathcal B_j,\label{F:drift}
\end{equation}
and the operator-norm difference between the recentered conditional covariance and $\mathcal B_j$
is at most $Ct\beta\operatorname{tr}\mathcal B_j$. Let $T$ be the vector martingale centered
under $Q$; then on the good set
\[
\|\langle T\rangle-\Sigma\|\le v+Ct\beta,\qquad|(U+\eta_kt,S+t)-T|\le C(tv+t^2\beta).
\]
The drift $-\eta_kt$ of the first coordinate comes from the cross entry of $\mathcal B_j$ in
\eqref{F:drift}; knowing only that the scalar variance is close to $1$ does not give this drift,
so the $EH$ bracket must be controlled simultaneously.

\emph{Step 3 (quantitative two-dimensional rectangle CLT).} For a centered two-dimensional
martingale $T$ with increment bound $b$, deterministic trace upper bound, and matrix bracket
within $w$ of $\Sigma$ with probability at least $1-q'$, set $\delta=b+w+q'$. For each unit
direction $\theta$, standardize $\theta\cdot T$ by $(\theta^\top\Sigma\theta)^{1/2}$ and stop when
the directional bracket first crosses $1+Cw$; the stopping changes only samples of probability at
most $q'$. A one-step Taylor expansion of the compensating characteristic function gives,
direction by direction, for $r=|u|\le\delta^{-1/16}$,
\[
|\varphi_T(u)-e^{-u^\top\Sigma u/2}|\le C\{br^3+b^2r^4+(w+b^2+q')r^2+q'\}\le C\delta(1+r^4),
\]
with constants uniform in the direction. Here only the characteristic function of the original
$T$ is estimated direction by direction; no common stopping vector is constructed.

Add an independent $N(0,h^2I_2)$ to both sides, with $h=\delta^{1/32}$. Fourier inversion gives
the supremum of the smoothed density difference: the integral over the low-frequency ball
$|u|\le\delta^{-1/16}$ is at most $C\delta^{5/8}$, and the high-frequency part is controlled by
$e^{-h^2|u|^2/2}$, which is superpolynomially small relative to $\delta$. The second moments of
both sides are uniformly bounded; restricting the lower rectangle to $[-R,R]^2$ with
$R=\delta^{-1/32}$, the outside probability is at most $C/R^2$ and the integral of the density
error inside is at most $CR^2\delta^{5/8}$, so the smoothed rectangle error is at most
$C\delta^{1/16}$. Finally remove the smoothing: take the coordinate tolerance
$a=h\sqrt{C\log(1/\delta)}$, so the probability that the noise exceeds the tolerance is less than
$\delta^2$; the marginal density of the smoothing Gaussian is uniformly bounded, and the
probability of the boundary strip of the rectangle is $O(a)$. Sandwiching the original rectangle
between the $\pm a$ shifts of the smoothed rectangle gives
\begin{equation}
d_{\rm rect}(T,N(0,\Sigma))\le C\delta^{1/64}.\label{F:rect}
\end{equation}
This conclusion also covers unbounded endpoints, that is, the marginal distributions.

Applying \eqref{F:rect} to the $T$ of Step 2 (with $b=2\beta$, $w=v+Ct\beta$), controlling the
shift $C(tv+t^2\beta)$ by the anticoncentration of the Gaussian marginals, and adding the total
variation error of Step 1 gives \eqref{F:rate}.
\end{proof}

\begin{remark}
The rate in \eqref{F:rate} is weak; it is used only to guarantee $t\epsilon_M\to0$, not to
predict the error at finite $M$.
\end{remark}

\subsubsection{Macro conditional CDF with all-positive signs}

\begin{proposition}[All-positive macro conditional law]\label{F:prop:cdf}
Under \eqref{F:range}, with the cutoff of Lemma~\ref{F:lem:cutoff}, the Plain model satisfies
\begin{equation}
\sup_x\left|\Pr(X_+\le x\mid \mathcal A)-\frac{\Phi_2(x,-t_M;\eta_k)}{\Phi(-t_M)}\right|\to0.\label{F:cdf}
\end{equation}
Equivalently, $X_++\eta_kt_M$ tends to $N(0,1-\eta_k^2)$ conditioned on $\mathcal A$.
\end{proposition}

\begin{proof}
\emph{Inverse tilting.} The stopped $S^*$ satisfies the scalar Laplace approximation
$\log\E e^{-tS^*}=t^2/2+o(1)$; take $t=b_M$. For any $y$, let
$\mathcal K_M(y)=\E_t[e^{t(S^*+t)}\mathbf1_{\{S^*+t\le0,\ Y^*+\eta_kt\le y\}}]$. A
Lebesgue--Stieltjes integration by parts writes $\mathcal K_M$ as an integral of the tilted joint
CDF, so Lemma~\ref{F:lem:tilt} gives $\sup_y|\mathcal K_M(y)-\mathcal K_G(y)|\le2\epsilon_M$,
where $\mathcal K_G$ corresponds to the target Gaussian. Also
$\mathcal K_G(\infty)=e^{t^2/2}\Phi(-t)\sim(\sqrt{2\pi}t)^{-1}$. The exact inverse-tilting
formula is
\[
\Pr(Y^*\le x,\ S^*\le-t)=\E e^{-tS^*}\,e^{-t^2}\,\mathcal K_M(x+\eta_kt),
\]
and controlling the numerator and denominator simultaneously by the difference between
$\mathcal K_M$ and $\mathcal K_G$ gives a conditional CDF error of at most
$O(\epsilon_M/\mathcal K_G(\infty))=O(t\epsilon_M)=o(1)$:
\[
\sup_x\left|\Pr(Y^*\le x\mid S^*\le-t)-\frac{\Phi_2(x,-t;\eta_k)}{\Phi(-t)}\right|\to0.
\]

\emph{Transfer to the macro event.} By input (C3) and the truncation and stopping bad sets of
this appendix, $\Pr(\mathcal A\triangle\{S^*\le-t\})=o(\Phi(-t))$. This probability bound
controls the joint subprobability error for all $\{Y_K\le x\}$ simultaneously; no energy-weighted
bound is needed, so the conditional CDF transfers. $Y_K$ and $Y^*$ differ by only a
superpolynomially small amount on the good set; the density of the selection Gaussian is uniformly
bounded, and given the score the conditional density of the energy is at most
$[2\pi(1-\eta_k^2)]^{-1/2}$, and so remains after mixing. Replace $t$ by $t_M$: the two
Gaussian-tail conditioning events are nested, so the total variation distance of the conditional
distributions is at most the relative tail difference $O(t_M|t-t_M|)=o(1)$. This gives the result
for $Y_K$ in the configuration model.

\emph{Fixed planting.} Fix $l$ special edges and construct the initial two-dimensional martingale
from $k(d-l)$ ordinary endpoints, then append the special-edge filtration of
Appendix~\ref{app:mda}. Replacing $\mu$ by $k(d-l)/M$ changes only an $O_l(M^{-1})$ error in
\eqref{F:bracketmean}. When the maximum ordinary degree is at most $L$, the $kl$ special
incidences change $K$ by at most $C_l(L+1)$ and $H$ by at most $C_l$; after normalization this is
$O_l(L/\sqrt M)$, which, placed in the perturbation of the energy coordinate and the bandwidth of
the macro event, still tends to zero after multiplying by $t$. Write
$G_M(x)=\Phi_2(x,-t_M;\eta_k)/\Phi(-t_M)$; then for each fixed $l$,
\begin{equation*}
\sup_x\left|\Pr\nolimits_l(Y_K\le x,\ \mathcal A)/p_M-G_M(x)\right|\to0.
\end{equation*}
This already contains the relative precision of the probability denominator, which an ordinary
finite-planting CLT cannot replace.

\emph{Plain transfer.} Let $N^{\rm bad}_M$ be the number of bad edges, with
$dp_M^{\rm bad}\to\Lambda_k$. Specifying $l$ bad edges is exactly the admissible special-edge
model, so for each fixed $l$,
\[
\sup_x\left|\E[\mathbf1_{\{Y_K\le x,\ \mathcal A\}}(N^{\rm bad}_M)_l]/p_M-\Lambda_k^lG_M(x)\right|\to0.
\]
Apply finite-order Bonferroni upper and lower bounds to the nonnegative indicator. First fix the
truncation order and let $M\to\infty$; the finitely many error terms are uniform in $x$. The
difference between the two-part sum and $e^{-\Lambda_k}$, multiplied by $G_M(x)\le1$, and then
the truncation order is increased. This gives
$\sup_x|\Pr(Y_K\le x,\mathcal A,N^{\rm bad}_M=0)/p_M-e^{-\Lambda_k}G_M(x)|\to0$; dividing by
$\Pr(\mathcal A,N^{\rm bad}_M=0)/p_M\to e^{-\Lambda_k}$ gives the conditional CDF in the Plain
model.

\emph{Convert back to $X_+$.} In the Plain model exactly
\begin{equation}
X_+=\frac{\sqrt2\mu\sqrt M}{k(1-k/M)s_M}Y_K+\frac{(k-1)\mu}{k(1-k/M)s_M},\label{F:affine}
\end{equation}
that is, $X_+=(1+O(M^{-1}))Y_K+O(M^{-1/2})$. Since $Y_K+\eta_kt_M$ is tight conditioned on
$\mathcal A$, the transformation error is $o_P(1)$, and the density upper bound above gives
\eqref{F:cdf}. The energy with all-negative signs is the same as with all-positive signs.
\end{proof}

\subsubsection{Conditional moments and uniform integrability}

Convergence in distribution does not give moment convergence after the center is shifted. This
section controls the tail directly under the macro conditioning.

\begin{proposition}[Macro conditional moments]\label{F:prop:moments}
Under the assumptions of Proposition~\ref{F:prop:cdf}, every fixed positive-integer moment and
absolute moment of $W^+_M=X_++\eta_kt_M$ conditioned on $\mathcal A$ converges to the
corresponding moment of $N(0,1-\eta_k^2)$. In particular,
\begin{equation}
\E[X_+\mid \mathcal A]=-\eta_kA(-t_M)+o(1),\qquad\Var(X_+\mid \mathcal A)=1-\eta_k^2+o(1).\label{F:moments}
\end{equation}
For any deterministic sign composition independent of the mapping,
$\E[X_s\mid \mathcal A]=\gamma\bigl[-\eta_kA(-t_M)+o(1)\bigr]$, where the $o(1)$ is the same
all-positive remainder, hence uniform in the composition and not requiring $\rho_M$ to converge.
\end{proposition}

\begin{proof}
\emph{Quantitative CDF input.} Work first in the configuration model. The polynomial rectangle
error of Lemma~\ref{F:lem:tilt}, the $t$ factor of the inverse tilting, the relative bandwidth
error $O(t(\delta_M+\epsilon_M))$, and the superpolynomially small bad set divided by $p_M$
together give
\begin{equation}
d_K\bigl(\mathcal L(Y_K\mid \mathcal A),\ \mathcal G_t\bigr)\le e_M=O(M^{-c}(\log M)^b),\label{F:dK}
\end{equation}
where $\mathcal G_t$ is the distribution of the first coordinate of a standard bivariate Gaussian
with correlation $\eta_k$, conditioned on the second coordinate being $\le-t$, and $c>0$,
$b<\infty$ are fixed. The Bonferroni bound has not yet been used here; later, transferring to the
Plain model needs only qualitative CDF convergence.

\emph{Tilted tail bound on the whole sample space.} Continue with the product tilting $Q$ of
\eqref{F:drift}. Under $Q$, $T_1=\sum_j(u_j-\E_Q[u_j\mid\mathcal F_{j-1}])$ is a centered
martingale. By \eqref{F:drift}, on the whole sample space
\[
\Bigl|\sum_j\E_Q[u_j\mid\mathcal F_{j-1}]\Bigr|\le t\sum_j|(\mathcal B_j)_{12}|+Ct^2\beta\sum_j\operatorname{tr}\mathcal B_j\le C_1t,
\]
using only $|(\mathcal B_j)_{12}|\le\operatorname{tr}\mathcal B_j/2$, the deterministic trace
bound, and $t\beta\to0$; it does not require the cross bracket to be close to $\eta_k$ on the bad
set. The increments of $T_1$ are at most $2\beta$ and the sum of predicted variances is at most
$C_2$, so the Bernstein--Freedman inequality gives
$Q(|T_1|>x)\le2\exp\{-x^2/[2(C_2+2\beta x/3)]\}$. Let $W^*=Y^*+\eta_kt$; then
$|W^*-T_1|\le C_3t$, so for $x\ge2C_3t$
\begin{equation}
Q(|W^*|>x)\le2\exp\{-cx^2/(1+\beta x)\}.\label{F:Qtail}
\end{equation}

\emph{Estimate unbounded moments without total variation.} The density of the true tilting with
respect to $Q$ is exactly $\prod_jm_j(-t)/\E e^{-tS^*}$, where
$m_j(-t)=\E(e^{-tv_j}\mid\mathcal F_{j-1})$. The deterministic trace bound and $t\beta\to0$ give
$\prod_jm_j\le e^{Ct^2}$; combined with $\E e^{-tS^*}\sim e^{t^2/2}$, this gives the
deterministic bound $dP_t/dQ\le e^{C_4t^2}\le M^{D_0}$ on the whole sample space, with $D_0$
fixed. This crude bound is used together with the martingale tail bound, whose order can be raised
arbitrarily; multiplying the small total variation distance by a growing maximum of the energy
does not yield a vanishing error.

\emph{Truncated tail.} Fix a positive integer $q$. The actual energy satisfies
$|Y_K|\le CM^{3/2}$, and the stopped martingale has at most $O(M)$ steps of size
$O(\log M/\sqrt M)$ each, so $|W^*|\le CM^2$ is a deterministic crude bound. Take
$R_M=A_q\sqrt{\log M}$ with $A_q$ large enough that $R_M\ge2C_3t$, and \eqref{F:Qtail} gives
$Q(|W^*|>R_M)\le2M^{-H_q}$ with $H_q>D_0+2q+2$. Outside $\mathcal B_M$, the event
$\mathcal A$ implies $S^*\le-t+C\delta_M$; for any nonnegative $f(W^*)$, the inverse tilting
gives
\[
\E[f(W^*);\mathcal A\setminus\mathcal B_M]\le\E e^{-tS^*}e^{-t^2}e^{Ct\delta_M}\E_tf(W^*)\le Cp_M\,t\,M^{D_0}\E_Qf(W^*),
\]
where the last step uses $\E e^{-tS^*}e^{-t^2}/p_M=O(t)$. Taking
$f=|W^*|^q\mathbf1_{\{|W^*|>R_M\}}$ gives
$\E[|W^*|^q\mathbf1_{\{|W^*|>R_M\}};\mathcal A\setminus\mathcal B_M]/p_M\le
CtM^{D_0+2q-H_q}=o(1)$. The contribution of the bad set is at most
$CM^{2q}\Pr(\mathcal B_M)/p_M=o(1)$. Hence
\begin{equation}
\E[|W_K|^q\mathbf1_{\{|W_K|>R_M\}}\mid \mathcal A]\to0,\qquad W_K=Y_K+\eta_kt.\label{F:tailcut}
\end{equation}
Since $R_M\to\infty$, \eqref{F:tailcut} is not yet uniform integrability by itself.

\emph{Between a fixed threshold and $R_M$.} Let $Z_0,G$ be independent standard normals; the
shifted-center variable of the selection Gaussian can be written as
$W^G_t=\eta_k(Z_0+t)+\sqrt{1-\eta_k^2}\,G$, conditioned on $Z_0\le-t$. For $t\ge1$ and $x\ge0$,
\[
\Pr(-(Z_0+t)>x\mid Z_0\le-t)=\frac{\Phi(-(t+x))}{\Phi(-t)}\le e^{-tx-x^2/2},
\]
which follows by shifting the integration variable of the numerator by $x$. Hence all moments of
the overshoot tend to zero, and the moments of $W^G_t$ tend to the corresponding moments of
$N(0,1-\eta_k^2)$. For the $q$-th moment use $\mathrm{clip}(w,-R_M,R_M)^q$, and for the absolute
moment use $\min(|w|,R_M)^q$; their total variation is at most $2R_M^q$, and Stieltjes integration
by parts together with \eqref{F:dK} gives that the truncated expectations differ by at most
$2R_M^qe_M=o(1)$. Removing the truncation via \eqref{F:tailcut} and the uniform tail control of
the Gaussian gives the convergence of all moments in the configuration model. Taking the
conclusion for order $q+1$ gives the uniform integrability
\begin{equation}
\sup_{M\text{ sufficiently large}}\E[|W_K|^q\mathbf1_{\{|W_K|>R\}}\mid \mathcal A]\le C_q/R.\label{F:UI}
\end{equation}

\emph{Plain transfer.} The scalar relative Bonferroni already gives
$\Pr_{\rm conf}(\mathcal A,N^{\rm bad}_M=0)/\Pr_{\rm conf}(\mathcal A)\to e^{-\Lambda_k}>0$.
For any nonnegative $f$,
\[
\E_{\rm Plain}[f(Y_K)\mid \mathcal A]=\frac{\E_{\rm conf}[f(Y_K)\mathbf1_{\{N^{\rm bad}_M=0\}}\mid \mathcal A]}{\Pr_{\rm conf}(N^{\rm bad}_M=0\mid \mathcal A)}\le C\,\E_{\rm conf}[f(Y_K)\mid \mathcal A],
\]
so \eqref{F:UI} transfers directly to the Plain model. This inequality transfers only the tail
control; the limiting distribution comes from the Plain CDF transfer of
Proposition~\ref{F:prop:cdf}. Together they give all moment limits in the Plain model. Finally
convert back to $X_+$ via \eqref{F:affine}:
$X_++\eta_kt_M=(1+O(M^{-1}))(Y_K+\eta_kt)+O(M^{-1/2})+\eta_k(t_M-t)+O(t/M)$, where the last
several terms tend to zero.

\emph{Mean and variance.} The first- and second-moment conclusions give
$\E[X_+\mid \mathcal A]=-\eta_kt_M+o(1)$ and
$\Var(X_+\mid \mathcal A)=1-\eta_k^2+o(1)$. The Mills ratio satisfies $A(-t)-t=O(1/t)$, so the
mean can also be written in the form \eqref{F:moments}. This is only a rewriting; the remainder is
still $o(1)$, and it does not identify the coefficient of order $1/t_M$.

\emph{Mean under an arbitrary composition.} The event $\mathcal A$ depends only on the unsigned
mapping and is invariant under permutations of the element labels, so the exact projection
\eqref{eq:projection} gives $\E[X_s\mid \mathcal A]=\gamma\E[X_+\mid \mathcal A]$. This step does
not require a sign CLT under a rare-event conditioning.
\end{proof}

\subsubsection{Conditional CDF and variance under an arbitrary sign composition, including
$|\rho|=1$}

\begin{proposition}[Macro conditional law under an arbitrary composition]\label{F:prop:sign}
Under the assumptions of Proposition~\ref{F:prop:cdf}, let the signs be an arbitrary deterministic
composition, with $\rho_M$ not required to converge. Let
$G_{M,s}(x)=\Phi_2(x,-t_M;\eta_k\rho_M^2)/\Phi(-t_M)$; then
\begin{equation}
\sup_x|\Pr(X_s\le x\mid \mathcal A)-G_{M,s}(x)|\to0,\qquad\Var(X_s\mid \mathcal A)=1-\eta_k^2\rho_M^4+o(1).\label{F:sign}
\end{equation}
\end{proposition}

The proof of this proposition is given after Lemma~\ref{F:lem:graph}--\ref{F:lem:endpoint}. The
idea is to avoid redoing a signed tilting for each sign composition. First generate an unsigned
Plain hypergraph, then assign the given numbers of positive and negative signs uniformly; this has
the same joint law of $(X_s,\mathcal A)$ as the original model, and $\mathcal A$ selects only the
graph. Hence two facts suffice: the required graph properties still hold with probability tending
to one under $\mathcal A$; and on these deterministic good graphs the sign CLT holds uniformly.
The graph properties used in Theorem~\ref{thm:sign} hold only to $o_P(1)$ accuracy and cannot be
divided directly by the rare $\Pr(\mathcal A)$. Below, the exceptional probabilities are
strengthened to be superpolynomially small.

Use the overlap graph notation of Appendix~\ref{sec:signproof}: $m_{ij}=|S_i\cap S_j|$,
$\Omega=\sum_{i<j}m_{ij}$, $\Omega_2=\sum_{i<j}m_{ij}^2$, $\ell_i=\sum_{j\ne i}m_{ij}$,
$\Omega_\ell=\sum_i\ell_i^2$, and $\bar\ell=2\Omega/d$.

\begin{lemma}[Superpolynomial concentration of the overlap graph statistics]\label{F:lem:graph}
There exists a deterministic set $\mathcal H_M$ of graphs such that
$\Pr(\mathcal G\notin\mathcal H_M)$ is superpolynomially small and all graphs in $\mathcal H_M$
uniformly satisfy
\begin{equation}
\frac\Omega M=\frac{\mu_c^2}2+o(1),\quad\frac{\Omega_2}M=\frac{\mu_c^2}2+o(1),\quad\frac1M\sum_i(\ell_i-\bar\ell)^2=\mu_c^2+o(1),\quad\max_i\ell_i\le kL.\label{F:goodgraph}
\end{equation}
All the $o(1)$ terms can be taken as $O((\log M)^4/\sqrt M+|r_M|/\sqrt M)$. In particular
$\Pr(\mathcal G\notin\mathcal H_M\mid \mathcal A)\to0$, and this exceptional probability
multiplied by any fixed power of $M$ and divided by $\Pr(\mathcal A)$ still tends to zero.
\end{lemma}

\begin{proof}
Let $L=\lceil\log M\rceil$ and $\mathcal D_L=\{\max_vD_v\le L\}$. Since
$D_v\sim\mathrm{Bin}(d,k/M)$, $\Pr(\mathcal D_L^c)\le M(eC/L)^L$, which is superpolynomially
small.

\emph{Lipschitz bound.} For two graphs in $\mathcal D_L$, replace the edges of differing labels
one by one by the corresponding edges of the other graph. The incidence sets of all intermediate
graphs are contained in the union of the two endpoint graphs, so the vertex degrees do not exceed
$2L$. Replacing one edge in such an intermediate graph affects only the $O(L)$ edge labels that
overlap the old or the new edge, and each $m_{ij}\le k$, so both $\Omega$ and $\Omega_2$ change
by $O(L)$. For the replaced edge $\ell_i=O(L)$, whose square changes by $O(L^2)$; for the
remaining $O(L)$ affected labels, $\ell_j$ changes by $O(1)$, so each square changes by $O(L)$.
Hence the total change of $\Omega_\ell$ is $O(L^2)$. Thus all three are $CL^2$-Lipschitz in the
Hamming distance of edge labels on $\mathcal D_L$. Here the global Lipschitz bound is proved along
a path in the union of two good graphs; it is not assumed that a single-step path stays within
$\mathcal D_L$.

\emph{Extension and concentration.} For each statistic $f$, take the McShane extension
$\widetilde f(\mathcal G)=\inf_{\mathcal G'\in\mathcal D_L}\{f(\mathcal G')+CL^2d_H(\mathcal G,\mathcal G')\}$
and truncate it to the deterministic nonnegative polynomial range of $f$. The extension equals
$f$ on the good set and remains $CL^2$-Lipschitz; the bad set has superpolynomially small
probability, and both $f$ and the truncated extension are at most $O(M^3)$, so
$|\E f-\E\widetilde f|$ is superpolynomially small. Applying McDiarmid's inequality to independent
edges,
\[
\Pr\bigl(|f-\E f|>C\sqrt M(\log M)^4\bigr)\le\Pr(\mathcal D_L^c)+2e^{-c(\log M)^4},\qquad f=\Omega,\Omega_2,\Omega_\ell.
\]

\emph{Deterministic means.} For the overlap $H$ of two independent uniform $k$-subsets, write
$h_1=\E H=k^2/M$ and $h_2=\E H^2=k^2/M+k^2(k-1)^2/[M(M-1)]$. Exactly,
\[
\E\Omega=\tbinom d2h_1,\qquad\E\Omega_2=\tbinom d2h_2,\qquad\E\Omega_\ell=d\{(d-1)h_2+(d-1)(d-2)h_1^2\}.
\]
The last identity holds because, given $S_i$, the $m_{ij}$ are independent and their conditional
first and second moments do not depend on $S_i$. Taking the good set as the intersection of
$\mathcal D_L$ with the three deviation threshold events gives \eqref{F:goodgraph}; the third
expression is computed from $\Omega_\ell-4\Omega^2/d$, in which the deterministic constants of the
two terms cancel. By input (C4), $\Pr(\mathcal A)\sim p_M$ is no smaller than a fixed negative
power divided by a polylogarithm, so the conditional exceptional probability tends to zero.
\end{proof}

\begin{lemma}[Uniform sign CLT on the good graphs]\label{F:lem:unifsign}
Make the exact decomposition $X_s=\gamma X_++V$ with $\E_s[V\mid\mathcal G]=0$. If
$\rho_M\to\rho$ and $|\rho|<1$, then for each fixed $u$,
\begin{equation}
\sup_{\mathcal G\in\mathcal H_M}\Bigl|\E_s[e^{iuV}\mid\mathcal G]-e^{-(1-\rho^4)u^2/2}\Bigr|\to0.\label{F:unifsign}
\end{equation}
\end{lemma}

\begin{proof}
The sign arguments of Appendix~\ref{sec:signproof} and Appendix~\ref{app:technical} are arguments
for deterministic graph sequences, using only \eqref{F:goodgraph}, $\Omega=O(M)$,
$\Omega_\ell=O(M)$, and a fixed distance of the sign proportions from the endpoints. Concretely,
the within-group variance bound for a whole-batch flip uses only $\Omega_\ell=O(M)$; the variance
of a flipped subset and the expectation of its internal edges use only $\Omega_\ell$ and $\Omega$.
The main term has $O(M)$ local terms, the maximum degree of the dependency graph and the magnitude
of the terms are both $O(\log M)$, and the normalized connected cumulant of fixed order $r\ge3$ is
$O_r(M^{1-r/2}\,\mathrm{polylog}\,M)$; the variance of the main term is determined by
\eqref{F:goodgraph}. Hence for any sequence $\mathcal G\in\mathcal H_M$,
$\mathcal L_s(V\mid\mathcal G)\Rightarrow N(0,1-\rho^4)$. If \eqref{F:unifsign} failed, one could
extract along a subsequence a deterministic graph violating it; this sequence of graphs still
satisfies all the deterministic bounds, contradicting the conclusion above.
\end{proof}

\begin{lemma}[Graph-by-graph conditional variance and the endpoints]\label{F:lem:endpoint}
Uniformly over all feasible compositions, on $\mathcal H_M$
\begin{equation}
\Var_s(V\mid\mathcal G)=1-\rho_M^4+o(1).\label{F:pergraphvar}
\end{equation}
\end{lemma}

\begin{proof}
The unconditional $\E V^2$ cannot be computed first and then divided by the rare $\Pr(\mathcal
A)$, so the graph-by-graph exact variance is used instead. Let $a_4$ be the mean of the sign
products of four distinct labels. The finite identity
\[
\Var_s\Bigl(\sum_{i<j}m_{ij}s_is_j\Bigm|\mathcal G\Bigr)=\Omega_2(1-2\gamma+a_4)+\Omega_\ell(\gamma-a_4)+\Omega^2(a_4-\gamma^2).
\]
Setting $v=1-\rho_M^2$ and expanding this finite rational expression uniformly over all feasible
compositions gives
\[
v^2\Omega_2+\rho_M^2v(\Omega_\ell-4\Omega^2/d)+O(\Omega_2/d+\Omega_\ell/d+\Omega^2/d^2).
\]
On $\mathcal H_M$ we have $\Omega,\Omega_2,\Omega_\ell=O(M)$, so the remainder is $O(1)$;
multiplying by the normalization factor $4/[k^2(1-k/M)^2s_M^2]\sim2/(\mu_c^2M)$ gives
\eqref{F:pergraphvar}. This also gives the crude upper bound $C(v+M^{-1})$. When
$|\rho_M|\to1$, \eqref{F:pergraphvar} and the graph-by-graph Chebyshev inequality give $V\to0$,
uniformly over the good graphs; this covers arbitrary rates of approach to the endpoints.
\end{proof}

\begin{proof}[Proof of Proposition~\ref{F:prop:sign}]
Let $W^+_M=X_++\eta_kt_M$. Proposition~\ref{F:prop:cdf} gives
$W^+_M\mid \mathcal A\Rightarrow N(0,1-\eta_k^2)$. For any subsequence with $\rho_M\to\rho$ and
$|\rho|<1$, condition first on the graph: on the good graphs the inner error tends to zero
uniformly by \eqref{F:unifsign}, on the bad graphs the absolute value of the characteristic
function is at most one, and by Lemma~\ref{F:lem:graph} the probability of a bad graph under
$\mathcal A$ tends to zero. Hence
\[
\E[e^{iaW^+_M+ibV}\mid \mathcal A]-\E[e^{iaW^+_M}\mid \mathcal A]\,e^{-(1-\rho^4)b^2/2}\to0,
\]
that is, $W^+_M$ and $V$ are asymptotically independent under $\mathcal A$. For $|\rho|=1$ use
instead the conditional probability convergence of Lemma~\ref{F:lem:endpoint}.

Keep the actual shifted center:
\[
X_s+\eta_k\rho_M^2t_M=\gamma W^+_M+V+\eta_k(\rho_M^2-\gamma)t_M,
\]
where the last term is $O(t_M/d)\to0$. Here $\rho_M^2t_M$ is not replaced by $\rho^2t_M$ in the
absence of a rate. Therefore
$\mathcal L(X_s+\eta_k\rho_M^2t_M\mid \mathcal A)\Rightarrow N(0,1-\eta_k^2\rho^4)$. The
selection benchmark $G_{M,s}$ has the same limit under the same shifted center: it can be
represented as $c(Z+t_M)+\sqrt{1-c^2}\,G$, conditioned on $Z\le-t_M$, where $c=\eta_k\rho_M^2$,
and the Gaussian overshoot tends to zero. The two distributions tend to the same continuous
distribution, so their CDF distance tends to zero. When $\rho_M$ does not converge, take a
subsequence on which the CDF distance has a positive lower bound, then by compactness of $[-1,1]$
extract a sub-subsequence on which $\rho_M$ converges, contradicting the conclusion above. This
proves the first display of \eqref{F:sign}.

The variance follows from the law of total variance:
\[
\Var(X_s\mid \mathcal A)=\gamma^2\Var(X_+\mid \mathcal A)+\E[\Var_s(V\mid\mathcal G)\mid \mathcal A].
\]
On the good graphs use \eqref{F:pergraphvar}; on the bad graphs use the deterministic polynomial
upper bound of the energy and the superpolynomial accuracy of Lemma~\ref{F:lem:graph}. Substituting
$\Var(X_+\mid \mathcal A)\to1-\eta_k^2$ from Proposition~\ref{F:prop:moments} and noting
$\gamma^2-\rho_M^4=O(1/d)$ gives the second display of \eqref{F:sign}. The conditional mean still
uses the exact projection and cannot be replaced by CDF convergence.
\end{proof}

\begin{remark}[Balancing rate]
Fixing $\rho_M\to0$ in the critical window is enough to recover the standard normal. In the
moderate deviation regime, the macro-layer center $-\eta_k\rho_M^2t_M$ grows with $t_M$: if
$\rho_M^2t_M\to0$ (which an exactly balanced composition satisfies), then
$X_s\mid \mathcal A\Rightarrow N(0,1)$; if $\rho_M^2t_M\to b\in(0,\infty)$, then the limit is
$N(-\eta_kb,1)$. This concerns the macro layer. When the weight of the repeated-edge layer in $F$
tends to one, the failure-conditioned law still recovers normality.
\end{remark}

\subsubsection{The two-mechanism mixture}

The failure-conditioned distribution is exactly a two-layer weighted mixture, for two reasons: a
failure that is neither a repeated edge nor a macro core has probability
$O(M^{-b_k})=o(\lambda_M)$; and the probability that a macro failure and a repeated-edge failure
occur simultaneously is also $o(\lambda_M)$. Concretely, Lemma~\ref{F:lem:cutoff} gives the exact
set decomposition $F=\mathcal A\cup D\cup R$ together with the sizes of the parts. The planted
transfer of Appendix~\ref{app:lowload} gives
$\sup_x|\Pr(X_s\le x\mid D)-\Phi(x)|\to0$, which requires only that $d/M$ tends to a finite
positive number and holds equally as $\alpha_M\to\alpha_c$. Hence
\[
\Pr(X_s\le x,F)=\Pr(\mathcal A)\Pr(X_s\le x\mid \mathcal A)+\lambda_M\Phi(x)+o(\lambda_M),
\]
with the error uniform in $x$. Substituting input (C4) and Proposition~\ref{F:prop:sign} gives
the mixture display of Proposition~\ref{prop:mixture}.

\subsection{Load-axis coverage}\label{app:axis}
This subsection proves Proposition~\ref{prop:axis}. The proof of coverage is assembled by load
interval, with each interval relying on a different result: the critical window uses
Theorem~\ref{thm:upper}; the moderate deviation regime uses Proposition~\ref{prop:mixture}; the
fixed positive subcritical regime uses Proposition~\ref{prop:lowload}; when $d/M\to0$ and
$d/\sqrt M\to\infty$, it uses the small-core counting with the $d$-dependence retained, together
with the CLT after planting a repeated edge; the moving subcritical regime
($|r_M|/\sqrt{\log M}\to\infty$) uses a sandwich via the monotonicity of the failure event under
edge-addition coupling; and the overload side, as well as $d/M\to\infty$, uses the unconditional
CLT along the load axis.

To distinguish it from the mixing rule of Section~\ref{sec:rules}, this appendix denotes the upper
bound with threshold $T_n$ in Proposition~\ref{prop:axis} by $U_{\rm mix0}$, and the
corresponding rule as the asymptotic mixing rule. The two differ in three respects: the asymptotic
mixing rule uses the unshifted coordinate $z_n$, uses $\lambda_n$ directly in the weight, and adds
no chi-square correction.

The mechanism of the counterexample can be seen from Proposition~\ref{prop:lowload}. On the sparse
side the failure almost always comes from a repeated edge; each repeated edge shifts
$\widehat d$ by $\pm2$, while $s_M$ is of order $d/\sqrt M$. When $d/\sqrt M\to\infty$ this shift
vanishes on the standardized scale; when $d=o(\sqrt M)$ it instead tends to infinity. Under
balanced signs the same-sign and opposite-sign repeated edges are each about half. The
opposite-sign half pushes the estimate too low, so the true value $d$ and all larger candidates
are rejected, and the probability that the upper bound misses the true value therefore tends to
$1/2$. For fixed $d=4$ with two positive and two negative signs, four of the six pairs are of
opposite sign, so this probability tends to $2/3$. Here the failure probability tends to zero, so
the positive lower bound required by Corollary~\ref{cor:balanced} fails, and balanced signs cannot
guarantee selection invariance.

In the sparse-side counterexample, $c_n\equiv0$, so the asymptotic mixing rule degenerates to the
unconditional normal threshold. In this case the failure-conditioned graph tends to a repeated
edge planted on a random pair, and the energy equals exactly $(d+2\sigma)/(1-k/M)$, where
$\sigma$ is the sign product of the repeated pair and satisfies
$\Pr(\sigma=-1)=d/[2(d-1)]$. Thus $U_{\rm mix0}=d+2\sigma$. The acceptance set is nonempty, so
the fallback for an empty set also does not rescue the coverage.

The coverage of the asymptotic mixing rule depends only on whether the true candidate is rejected.
Let $\epsilon_M=\sup_x|\Pr(X\le x\mid F)-H_d(x)|$. Since $H_d$ is continuous, taking left limits
gives, for every real $q$, $|\Pr(X<q\mid F)-H_d(q)|\le\epsilon_M$; take $q=H_d^{-1}(u)$. So it
suffices to prove $\epsilon_M\to0$. By a subsequence argument it suffices to treat the following
cases separately.

Inside the critical window, $\lambda_M\to0$ and $H_d$ returns to the selection law. The moderate
deviation regime is handled by Proposition~\ref{prop:mixture}. The fixed subcritical regime has
$w_d\to0$ and is handled by Proposition~\ref{prop:lowload}.

When $r_M\to-\infty$, $|r_M|/\sqrt{\log M}\to\infty$, and $d/M\to\alpha_c$, take
$a=g_k^2C^2/2>k-2$ and $d'_M=\lfloor\alpha_cM-C\sqrt{M\log M}\rfloor$, so that eventually
$d_M\le d'_M$. Under edge-addition coupling the edge set of the core is monotone under edge
addition, so by Proposition~\ref{prop:mda},
$\Pr_{d_M}(\mathcal A)\le\Pr_{d'_M}(\mathcal A)=M^{-a+o(1)}=o(\lambda_M)$. Thus
$\Pr(F\setminus D)=o(\lambda_M)$, the conditional law is determined by the repeated-edge layer,
and $H_d\to\Phi$.

On the overload side and when $d/M\to\infty$, $\Pr(F)\to1$, and the conclusion follows from the
unconditional CLT along the load axis. This CLT holds for every sequence with
$d/\sqrt M\to\infty$. The reason is that the edge-revealing martingale of Appendix~\ref{app:lowload}
satisfies $\Var\mathcal V_d=O(d^4/M^3)$, while $s_M^4\approx4d^4/M^2$, so the variance of the
normalized bracket is $O(1/M)$; the sum of fourth moments satisfies
\begin{equation*}
\frac{\sum_j\E\Delta_j^4}{s_M^4}\le C_k\Bigl(\frac M{d^2}+\frac1d\Bigr)\to0,
\end{equation*}
so the martingale CLT gives $X\Rightarrow N(0,1)$. This conclusion does not require $d/M$ to be
bounded, and it is uniform in the signs. When a repeated edge is planted on a random pair, the
deterministic term $2\sigma$ and the cross term, divided by $s_M$, have squared errors
$O(M/d^2)$ and $O(1/d)$ respectively, which also tend to zero.

\paragraph{The $d/M\to0$ conditional transfer.} Here the dependence on $d$ in the small-core
counting of Appendix~\ref{app:lowload} is retained. In this regime every possible core size
satisfies $e\le d=o(M)$, and
\[
B_e\le(C_kd/M)^e(e/M)^{\lceil(k-2)e/2\rceil}.
\]
Retaining the rounding for $e=3$ and summing the upper bound over $e\ge4$ gives
\[
\Pr(F\setminus D)\le C_k\left(\frac{d^3}{M^{b_k+3}}+\frac{d^4}{M^{2k}}\right),
\qquad b_k=\lceil3(k-2)/2\rceil\ge k-1.
\]
By $\lambda_M\asymp_k d^2/M^k$ and the second-factorial-moment bound of repeated pairs, the
random-pair repeated-planting distribution $Q$ satisfies
\[
\bigl\|P(\cdot\mid F)-Q\bigr\|_{\rm TV}
\le C_k\left(\lambda_M+dM^{k-b_k-3}+\frac{d^2}{M^k}\right)\longrightarrow0.
\]
Hence when $d/\sqrt M\to\infty$, the planted CLT above gives $X\mid F\Rightarrow N(0,1)$. On the
other hand, $z_M(d)=-(g_k\alpha_c+o(1))\sqrt M$, while $\lambda_M\ge1/\binom Mk$, so $w_d\to0$
and $H_d\to\Phi$ uniformly. This completes the branch with load tending to zero. For a
subsequence, first take $d/M$ converging in $[0,\infty]$; when the limit is $\alpha_c$, then take
the limit of $r_M$ and of $|r_M|/\sqrt{\log M}$. This covers all load sequences.

Now construct the sparse-side counterexample. When $b=0$, each candidate has $c_n=0$, and the
asymptotic mixing rule degenerates exactly to the unconditional normal threshold
$T_n=n+s(n)\Phi^{-1}(u)$. When $d=o(\sqrt M)$, the failure-conditioned graph tends to a repeated
edge planted on a random pair of labels; the probability that the remaining base edges are pairwise
disjoint is $1-O(d^2/M)$, in which case $\widehat d=(d+2\sigma)/(1-k/M)$. For $M$ sufficiently
large, $T_n$ is strictly increasing in $n\ge2$, and
$T_{d+2\sigma}<\widehat d<T_{d+2\sigma+2}$, so $U_{\rm mix0}=d+2\sigma$, hence
$\Pr(U_{\rm mix0}<d\mid F)\to\Pr(\sigma=-1)=d/[2(d-1)]$.
\subsection{The second sign change of the mean}\label{app:wm}
Finally consider the conditional mean. In the fixed subcritical regime $c_F$ is positive
(Proposition~\ref{prop:lowload}); in the critical window $c_F$ is negative for $k\ge5$
(Section~\ref{sec:flip}). So for $k\ge5$, as the load decreases, $c_F$ must change sign once
more. The weighted formula of this section uses the all-same-sign measurement: set all $s_i=+1$
and write $Y=\widehat d-d=s_MX_{\rm one}$.

\begin{proposition}[Weighted mean and the sign-change layer of the mean for $k\ge5$]\label{prop:wm}
In the moderate deviation range of Section~\ref{sec:loadaxis}, with the cutoff of
Lemma~\ref{F:lem:cutoff},
\begin{equation*}
\E[Y;F]=2\lambda_M-s_M\eta_k\phi(t_M)+o\bigl(\lambda_M+s_M\phi(t_M)\bigr).
\end{equation*}
Fix $k\ge5$ and let
\begin{equation*}
t_*^2=(2k-3)\log M-2\log\frac{\alpha_ck!\sqrt\pi}{\eta_k}.
\end{equation*}
For any fixed $v\in\mathbb R$, take integers $d_M$ with $t_M^2-t_*^2\to v$; then
$c_F\binom d2\to1-e^{-v/2}$. Hence between the two layers $t^2=t_*^2\pm\delta$, $c_F$ changes
sign at least once.
\end{proposition}
The first display of the proposition is proved at the end of this subsection, and the second part
follows from the first; the meaning of the two conclusions is explained first. For any
deterministic sign composition, $\E[\widehat d-d;F]$ is exactly the all-same-sign expectation
multiplied by the sign factor $\gamma=(d\rho_M^2-1)/(d-1)$, so both the main term and the
remainder are projected by this factor. The coefficient $c_F$ is extracted from the all-same-sign
measurement, so the above conclusion about the sign-change layer is unaffected.

Why is the sign-change layer deeper than the probability-equal-weight layer? The equal-weight point
compares the failure probabilities of the two mechanisms, namely $\Phi(-t)\approx\lambda_M$,
which by the Mills ratio gives $t^2\approx(2k-4)\log M-\log\log M$. The sign change of the mean
compares instead the contributions of the two mechanisms to the mean, namely
$s_M\phi(t)\approx\lambda_M$. The factor $A(-t)=\phi(t)/\Phi(-t)$ in the macro conditional mean
cancels the Mills factor exactly, so the sign-change layer lies at
$t^2\approx(2k-3)\log M$, deeper than the equal-weight layer, and without a $\log\log M$ term. On
the sign-change layer the weight of the macro failure is only $O(1/\sqrt{M\log M})$; but the
energy shift that a macro failure produces on the original scale is $O(\sqrt{M\log M})$, and the
product is still enough to cancel the $+2$ of the repeated-edge layer. Thus the mechanism that
dominates the failure probability is not the mechanism that dominates the conditional mean.
Figure~\ref{fig:loadaxis}(a) plots the position of the two layers for $k=5$ at finite size.

The two kinds of sign change must also be distinguished. The sign change of Section~\ref{sec:flip}
fixes the load and varies $k$; the sign change here fixes $k\ge5$ and moves along the load. They
are different phenomena.

The $o(1)$ error of the mixture CDF cannot be multiplied by a growing macro center, so the
conditional mean must be controlled separately on the unnormalized first moment.

Now prove the first display of Proposition~\ref{prop:wm}. Let $N$ be the number of repeated edge
pairs and $D=\{N\ge1\}$.

\begin{proof}
The exact density of the random-pair planting gives $\E[YN]=2\lambda_M$. Let
$R=F\setminus(\mathcal A\cup D)$; there is the graph-by-graph identity
\begin{equation}
Y\mathbf1_F-Y\mathbf1_{\mathcal A}-YN=Y\mathbf1_R-YN\mathbf1_{\mathcal A}-Y(N-1)\mathbf1_{D\setminus \mathcal A}.\label{F:wm-id}
\end{equation}
The absolute first moments of the three terms on the right are estimated one by one below.

\emph{Third-kind failures.} The moment estimate over covering events in the mean expansion of
Appendix~\ref{app:lowload} uses only $d\le AM$, so it holds uniformly on the present moving load:
$\E[|Y|;R]\le C\sqrt M\,M^{-b_k}=o(\lambda_M)$, since $b_k-(k-2)\ge1$. Here the core on $R$ does
not exceed $\varepsilon_0M$ edges, and no fixed subcritical gap is needed.

\emph{Multiplicity of repeated pairs.} The joint conditioning on any two prescribed repeated pairs
involves at most four special labels, and the internal, cross, and external decomposition of the
energy still gives $\E|Y|\le C\sqrt M$; also $\E[N(N-1)]\le\lambda_M^2$, so
$\E[|Y|(N-1);D]\le C\sqrt M\lambda_M^2=o(\lambda_M)$.

\emph{Intersection of the macro layer and the repeated layer.} Let $Q$ be the repeated-planting
distribution of a random pair, whose density is exactly $N/\lambda_M$. Then
\[
\E[|Y|N;\mathcal A]=\lambda_M\E_Q[|Y|;\mathcal A]\le\lambda_M\sqrt{\E_QY^2}\sqrt{Q(\mathcal A)}\le C\lambda_M\sqrt{Mp_M}.
\]
$\E_QY^2=O(M)$ comes from the energy decomposition after repeated planting and the external
unconditional variance of Appendix~\ref{app:lowload}. $Q(\mathcal A)\sim p_M$ comes from the
fixed-planting version of input (C4): take the two identical valid $k$-subsets as finite special
edges, then apply a relative Bonferroni transfer to the remaining ordinary edges. Let
$b=s_Mt_M/\lambda_M$; by $\sqrt{p}/(1+bp)\le1/(2\sqrt b)$,
\[
\frac{\lambda_M\sqrt{Mp_M}}{\lambda_M+s_Mp_Mt_M}\le\frac12\sqrt{\frac{M\lambda_M}{s_Mt_M}}=O\Bigl(M^{(5/2-k)/2}/\sqrt{t_M}\Bigr)=o(1).
\]
When $t_M\to\infty$, $p_MA(-t_M)=\phi(t_M)\approx p_Mt_M$, so the intersection term is
$o(\lambda_M+s_M\phi(t_M))$, and no refined weighted local response is needed.

\emph{Macro main term.} By input (C4) and Proposition~\ref{F:prop:moments},
$\E[Y;\mathcal A]=s_M\Pr(\mathcal A)[-\eta_kA(-t_M)+o(1)]=-s_M\eta_k\phi(t_M)+o(s_M\phi(t_M))$.
The probability needs only relative $o(1)$ accuracy, since the target remainder is
$o(\lambda_M+s_M\phi(t_M))$.

Substituting the three estimates and the main term into \eqref{F:wm-id} gives the first display of
Proposition~\ref{prop:wm}.
\end{proof}

Combining $\Pr(F)=(p_M+\lambda_M)(1+o(1))$ with $c_F\binom d2=\E[Y\mid F]/2$ gives
\begin{equation*}
c_F\binom d2=\frac{\lambda_M-s_M\eta_k\phi(t_M)/2}{p_M+\lambda_M}+o\Bigl(\frac{\lambda_M+s_M\phi(t_M)}{p_M+\lambda_M}\Bigr).
\end{equation*}
For $k\ge5$, $t_M^2-t_*^2\to v$ gives
$s_M\eta_k\phi(t_M)/(2\lambda_M)\to e^{-v/2}$ and $p_M/\lambda_M\to0$, hence
$c_F\binom d2\to1-e^{-v/2}$.
